\clearpage
\section[Rencana Implementasi]{}
\sectionopener{cgreen}{Peta Jalan}{Rencana Implementasi}{Atlas tidak dimulai dari nol pasca-kompetisi: pondasi produksi sudah berjalan, rencana ke depan adalah memperluas jangkauan dan memperkuat keandalannya.}

\diagrambox{assets/diagrams/roadmap.pdf}{1}

\vspace{1em}
\subsection{Rincian setiap tahap}

\begin{enumerate}[leftmargin=1.4em]
  \item \textbf{Sekarang, MVP produksi.} Seluruh modul (PMO, chat, suara, catatan/whiteboard real-time, Godmode) berjalan dan dipakai harian oleh Laboratorium MGM; bersumber terbuka di bawah AGPL-3.0.
  \item \textbf{Jangka pendek.} Perluasan pengujian bersama komunitas dan himpunan kampus lain, penguatan audit keamanan (terutama alur otentikasi multi-metode), serta perbaikan berdasarkan umpan balik pemakai nyata di luar MGM Laboratory.
  \item \textbf{Menengah.} Aplikasi mobile pendamping, integrasi kalender dan email eksternal, serta ekspor/impor data untuk migrasi dari alat lain (Jira, Trello, Notion).
  \item \textbf{Panjang.} Sistem plugin/integrasi pihak ketiga, dan dukungan federasi multi-instans agar organisasi besar dapat menghubungkan beberapa deployment Atlas tanpa kehilangan kendali data masing-masing.
\end{enumerate}

\vspace{0.6em}
\begin{calloutbox}{cgreen}
{\bricolagesb\fontsize{11pt}{14pt}\selectfont\color{cink} Mengapa peta jalan ini realistis}\\[0.3em]
{\fontsize{9.6pt}{13pt}\selectfont\color{cink2} Setiap tahap dibangun di atas kode yang sudah berjalan produksi, bukan proposal fitur yang belum tersentuh. Tim tidak perlu meyakinkan pemakai baru bahwa Atlas ``akan'' berfungsi, sebab 65+ pemakai nyata sudah membuktikannya setiap hari.}
\end{calloutbox}

\clearpage
